% Created 2025-11-11 Tue 18:31
% Intended LaTeX compiler: pdflatex
\documentclass[11pt]{article}
\usepackage[utf8]{inputenc}
\usepackage[T1]{fontenc}
\usepackage{graphicx}
\usepackage{longtable}
\usepackage{wrapfig}
\usepackage{rotating}
\usepackage[normalem]{ulem}
\usepackage{amsmath}
\usepackage{amssymb}
\usepackage{capt-of}
\usepackage{hyperref}
\author{Amitav Krishna}
\date{\today}
\title{}
\hypersetup{
 pdfauthor={Amitav Krishna},
 pdftitle={},
 pdfkeywords={},
 pdfsubject={},
 pdfcreator={},
 pdflang={English}}
\begin{document}

\tableofcontents

\hyperref[sec:org48a6042]{\uline{Introduction 1}}

\hyperref[sec:orge5f032e]{\uline{Methods 1}}

\begin{quote}
\hyperref[sec:org826e1f8]{\uline{2.1
Quantum Circuit Dataset Generation and Density Matrix Extraction 1}}

\hyperref[sec:orged7fdbf]{\uline{2.2 Noise Modelling 2}}

\hyperref[sec:org993ec26]{\uline{2.3 Transformer architecture 2}}

\hyperref[sec:org18a288f]{\uline{2.4 Physics-informed loss 2}}

\hyperref[sec:orgfc9c916]{\uline{2.5 Training 3}}
\end{quote}

\hyperref[sec:org1863f1c]{\uline{Analysis 3}}

\hyperref[sec:org69a40e4]{\uline{Related Implementations 4}}

\hyperref[sec:orga20fc04]{\uline{Conclusion 5}}
\section{Introduction}
\label{sec:org48a6042}
Quantum computation and sensing has the potential to revolutionize
industries from cryptography to medicine. Right now, the effectiveness
of these systems are severely handicapped by the presence of persistent,
inherent noise. We propose a physics-informed, transformer-based
denoising autoencoder for quantum state noise mitigation. We believe
that the properties of non-locality inherent to transformers offer a
superior reflection of the non-local nature of quantum entanglement,
compared to the inherently local perspective imposed by traditional CNN
architectures.
\section{Methods}
\label{sec:orge5f032e}
\subsection{Quantum Circuit Dataset Generation and Density Matrix Extraction}
\label{sec:org826e1f8}
In order to aid supervised learning, we fabricate a large scale,
synthetic dataset with known "ground truth." We use a quantum
programming library such as Cirq to generate 10,000 random 5-qubit
random quantum circuits. The circuits are implemented in a controlled
manner such that we cover different noise regimes: there are 500
instances for each combination of five types of noises, and four levels
of intensity (5 \textbf{4} 500 = 10,000 circuits). In addition, we use an 80/20
train-test split prior to training to ensure unbiased testing.
\subsection{Noise Modelling}
\label{sec:orged7fdbf}
The resulting noisy circuits are simulated with Cirq's
DensityMatrixSimulator, and are the inputs that will be given to the
network. The types of noise applied are as followed:

\begin{itemize}
\item \textbf{Bitflip:} Introduces random Pauli-X operations on qubits, flipping
the computational basis state
\href{https://www.codecogs.com/eqnedit.php?latex=\%5C\%7C0\%5Crangle\%20\%5Cleftrightarrow\%20\%5C\%7C1\%5Crangle\#0}{{[}[file:media/image5.png}]]
with a given probability.

\item \textbf{Depolarizing:} With a given probability, replaces the qubit's state
by the maximally mixed state, reducing its purity uniformly across all
bases.

\item \textbf{Amplitude damping:} Models energy relaxation, where the excited state
\href{https://www.codecogs.com/eqnedit.php?latex=\%7C1\%5Crangle\#0}{{[}[file:media/image4.png}]]
decays to the ground state
\href{https://www.codecogs.com/eqnedit.php?latex=\%7C0\%5Crangle\#0}{{[}[file:media/image1.png}]],
mimicking spontaneous emission or photon loss.

\item \textbf{Phase damping:} Represents loss of phase coherence between states
without any exchange of energy. Populations remain unchanged while
off-diagonal terms decay.

\item \textbf{Mixed:} A composite channel that applies a random combination of the
above noise mechanisms with equal likelihood, approximating general
environmental decoherence.
\end{itemize}
\subsection{Transformer architecture}
\label{sec:org993ec26}
We used an autoencoding transformer architecture trained to map noisy
density matrices onto the clean ones. The architecture is selected for
its Transformer-like advantage in capturing long-range, non-local
correlations -- a necessary condition to adequately reconstruct
entanglement and coherence in mult

i-qubit states. The workflow is as follows: The input density matrix is
projected into a sequence of patch tokens. This sequence is transformed
within the encoder based Transformer network using self-attention layers
to capture global noise patterns. The output is then reconstructed via a
lightweight decoder to yield the denoised state ,
\subsection{Physics-informed loss}
\label{sec:org18a288f}
Following the physics-informed neural network (PINN) paradigm

We design the training objective to embed quantum mechanical structure
directly into the optimization process. The total loss combines a
data-driven reconstruction term, mixing mean-squared error, frobenius
fidelity, and physical regularizers that enforce Hermiticity, unit
trace, and positive semidefiniteness of the predicted density matrices.
This formulation constrains the network outputs to the manifold of valid
quantum states, ensuring that denoised predictions remain consistent
with the postulates of quantum mechanics.
\subsection{Training}
\label{sec:orgfc9c916}
The model was trained for 100 epochs with a batch size of 16. The
training process employed the Adam optimizer with a learning rate of
3e-4. The training was conducted on an NVIDIA T4 GPU, with one epoch
taking approximately 22 seconds. Model checkpoints were saved every
epoch, and no learning rate scheduling or early stopping was employed.

The Transformer autoencoder uses a patch embedding size of 2×2, an
embedding dimension of 128, 6 encoder blocks, and 8 attention heads per
block. Each block includes a 2× expansion MLP and standard pre-norm
residual connections. The decoder consists of a single transposed
convolution layer projecting back to a 2-channel density matrix
representation.

To improve robustness across heterogeneous noise regimes, we adopoted an
\textbf{adversarial reweighting scheme} during training. Specifically, losses
were grouped by noise intensity levels (0.05, 0.10, 0.15, 0.20), and
adaptive group weights were updated via an exponentiated-gradient rule:

\begin{center}
\includegraphics[width=.9\linewidth]{media/image2.png}
\end{center}

Where $$L_i$$ is the mean loss for group $$i$$ and $$\eta = 0.05$$
controls the adversarial learning rate. This mechanism dynamically
upweights samples from higher noise distributions, encouraging the model
to prioritize difficult noise conditions and improving generalization
across the full noise spectrum.
\section{Analysis}
\label{sec:org1863f1c}
\subsection{Performance}
\label{sec:org8dc812e}
After training, the model achieved an average fidelity improvement of
\textbf{0.133} across all noise conditions. However, this gain was not
uniformly distributed. As shown in Figure 1, performance deteriorates
sharply at higher noise levels due to inputs that lack discernible
quantum structure, resulting in low-information or effectively random
reconstructions. The transformer contains over 800,000 parameters, while
the dataset contains a relatively small number of circuits in
comparison, boasting just 10,000. Transformers typically require
substantially larger datasets to learn robust, global noise-correlation
structures, and this mismatch likely contributes to the weaker
generalization observed at high noise levels compared to convolutional
baselines[Kendre, 2025]. Furthermore, with limited data, the
\textbf{physics-informed regularizers} dominate the loss, biasing the model
toward producing physically valid, but overly conservative, density
matrices rather than closely reconstructing the true clean states. In
effect, this over-regularization can degrade the denoising process
itself, yieliding negative fidelity deltas in cases where the network
overcorrects physically valid but noisy states\begin{center}
\includegraphics[width=.9\linewidth]{media/image3.png}
\end{center}
\subsection{Scalability}
\label{sec:orgcd640de}
These limitations are largely attributable to the data-model scale
mismatch rather than the architecture itself. Transformers are known to
exhibit strong data-scaling laws: as both model capacity and dataset
size increase, their ability to capture long-range dependencies improves
predictably[Kaplan et al., 2020; Hoffmann et al. 2022]. In this context,
expanding the dataset by several orders of magnitude would allow the
network to better model complex, non-local noise correlations that
convolutional architectures inherently miss[Vaswani et al., 2017;
Dosovitskiy et al., 201]. Furthermore, larger-scale training would
rebalance the optimization landscape, reducing the dominance of the
physics-informed regularizers and enabling the model to prioritize
reconstruction fidelity without sacrificing physical validity[Norambeuna
et al., 2023; Raissi et al., 2019]. This suggests that the transformer
framework remains a promising foundation for quantum-state denoising,
provided it is trained in a sufficiently data-rich regime.
\section{Related Implementations}
\label{sec:org69a40e4}
\begin{adjustbox}{max width=\textwidth}
\small
\begin{center}
\begin{tabular}{llll}
Feature & Machine learning for Quantum Noise Reduction [Kendre, 2025] & Physics-Informed Transformers for Quantum Noise Reduction & Learning to Restore Heisenberg Limit in Noisy Quantum Sensing via Quantum Digital Twin [Xu et al., 2025]\\
\hline
Core Architecture & Convolutional Neural Network (CNN) autoencoder & Transformer autoencoder & Quantum Digital Twin Sensing (QDTS) Protocol\\
Loss Function & Used a supervised learning approach with a fidelity-aware composite loss function & Uses a Physics-Informed Loss enforcing Hermiticity, unit trace, positive semidefiniteness & MSE for observables + RL agent reward = Classical Fisher Information\\
Data and Scope & 10,000 noisy 5-qubit density matrices & Similar: 10,000 random 5-qubit quantum circuits & Continuous environmental measurement (photon counting, homodyne) across NISQ platforms\\
Key Limitation & CNNs miss non-local entanglement structure & Self-attention captures long-range entanglement and coherence & Needs continuous monitoring + large data; no experimental validation yet\\
Performance & ΔFidelity = 0.476 & ΔFidelity = 0.133 & Restores or surpasses Heisenberg Limit in simulation\\
\end{tabular}
\end{center}
\end{adjustbox}
\section{Conclusion}
\label{sec:orga20fc04}
In conclusion, this work introduces the Quantum-Fidelity
Physics-Informed Denoising Transformer (QFPIDT), a deep learning
framework designed to mitigate quantum noise through a fidelity-aware
optimization process. While initial experiments revealed high variance
and limited fidelity under mixed noise conditions, these findings
highlighted the intrinsic complexity of quantum noise modeling and
motivated the transition from convolutional to Transformer-based
denoising architectures. Leveraging global self-attention, the
Transformer model demonstrates enhanced capacity to restore coherence by
capturing non-local dependencies within quantum states.

Future work will involve scaling the training corpus to millions of
simulated quantum circuits. Such large-scale data is expected to better
exploit the data-hungry nature of attention mechanisms, enabling the
Transformer to generalize across diverse noise regimes. Moreover, a
proposed follow-up study will compare a standard Transformer with a
\textbf{physics-informed Transformer}---one constrained by quantum-mechanical
priors---to assess how embedding physical laws within the model
encourages \emph{truth-seeking} representations over purely statistical
correctness. This direction holds promise for developing denoising
architectures that bridge empirical learning with fundamental quantum
principles, advancing the viability of QFDA in real-world biosensing
applications.
\end{document}
